\documentclass[10pt,a4paper]{article}

\usepackage[margin=1.35cm]{geometry}
\usepackage{fontspec}
\setmainfont{DejaVu Sans}
\usepackage[english,russian]{babel}
\usepackage{xcolor}
\usepackage{hyperref}
\usepackage{enumitem}
\usepackage{titlesec}
\usepackage{tabularx}
\usepackage{microtype}

\pagestyle{empty}
\setlength{\parindent}{0pt}
\setlength{\parskip}{0pt}
\setlength{\emergencystretch}{2em}
\linespread{1.04}
\setlist[itemize]{
  leftmargin=1.35em,
  itemsep=2pt,
  topsep=4pt,
  partopsep=0pt,
  parsep=0pt
}

\definecolor{accent}{HTML}{1F4E79}
\definecolor{muted}{HTML}{555555}
\hypersetup{
  colorlinks=true,
  urlcolor=accent,
  linkcolor=accent,
  pdfauthor={Дмитрий Гудов},
  pdftitle={Резюме — Дмитрий Гудов}
}

\titleformat{\section}
  {\large\bfseries\color{accent}}
  {}{0pt}{}
  [\color{accent}\titlerule]
\titlespacing*{\section}{0pt}{10pt}{5pt}

\newcommand{\entry}[3]{%
  \noindent\begin{tabularx}{\textwidth}{@{}Xr@{}}
    \textbf{#1} & \textcolor{muted}{#2}
  \end{tabularx}\par
  \vspace{1pt}
  #3\par\vspace{3pt}
}
\newcommand{\project}[2]{\textbf{#1} — #2\par\vspace{3pt}}

\begin{document}
\raggedright
\begin{minipage}{\textwidth}
{\LARGE\bfseries Дмитрий Гудов}\par
\vspace{3pt}
{\large\color{accent} DevOps / Linux-инженер · системный и backend-разработчик}\par
\vspace{6pt}
\href{mailto:gudovdo@my.msu.ru}{gudovdo@my.msu.ru}
\quad·\quad
\href{https://github.com/kr0sh512}{github.com/kr0sh512}
\quad·\quad
\href{https://t.me/kr0sh_512}{t.me/kr0sh\_512}
\end{minipage}\par

\section{Профессиональный профиль}
DevOps-ориентированный разработчик и студент ВМК МГУ. Разрабатываю и сопровождаю Linux-инфраструктуру, автоматизирую повторяемые операции и создаю системные и backend-инструменты. Основные направления: ALT Linux и внутреннее устройство RPM, Kubernetes, контейнеризация, CI/CD, Python и C.

\section{Опыт}
\entry{BaseALT}{по настоящее время}{Инструменты для Linux/RPM и системные исследования: механизмы зависимостей ALT Linux, производительность и совместимость RPM.}

\entry{Яндекс — стажёр, команда DevTools автономного транспорта}{осень 2025}{Разработка в команде инструментов для разработчиков.}

\section{Образование}
\entry{Московский государственный университет имени М. В. Ломоносова}{обучаюсь}{Факультет вычислительной математики и кибернетики, Computer Science.}

\section{Ключевые навыки}
\begin{itemize}
  \item \textbf{DevOps и инфраструктура:} Linux, Docker/OCI, Kubernetes, K3s, RKE2, Helm, Ansible, Traefik, systemd, SSH, Let's Encrypt.
  \item \textbf{CI/CD:} GitHub Actions, GHCR, мультиархитектурные сборки контейнеров, автоматизированное тестирование, упаковка и публикация.
  \item \textbf{Разработка:} Python, C, C++, TypeScript, Shell, SQL; FastAPI, SQLAlchemy, Pydantic, React, Vite.
  \item \textbf{Данные и качество:} PostgreSQL, SQLite, MongoDB; pytest, mypy, Black, санитайзеры, дифференциальное тестирование, бенчмаркинг.
\end{itemize}

\section{Избранные проекты}
\project{\href{https://github.com/kr0sh512/alt-rpm-set-version}{ALT RPM Set Version Research}}{исследование кодированных множеств символов и сравнения зависимостей RPM: C-реализации, хеширование, совместимость, проверки с санитайзерами и APT/RPM-бенчмарки.}
\project{\href{https://github.com/kr0sh512/nexus-sync}{nexus-sync}}{клиент-серверная система централизованного управления компьютерами и запуска разрешённых команд. FastAPI, SQLAlchemy, CLI, YAML-конфигурация, systemd, автономные сборки и CI.}
\project{Персональная инфраструктура}{автоматизация Linux-серверов через Ansible, эксплуатация многоузловых K3s/RKE2, контейнерные развёртывания, ingress/ServiceLB, автоматизация TLS и удалённая диагностика.}
\project{\href{https://github.com/kr0sh512/pomodoro}{Pomodoro}}{браузерный таймер на React, TypeScript и Vite; тестируется и публикуется GitHub Actions как непривилегированный мультиархитектурный образ в GHCR.}
\project{\href{https://github.com/kr0sh512/network-diagrams-tool}{Network Diagrams Tool}}{Python-инструмент для построения и визуализации сетевых топологий.}
\project{\href{https://github.com/kr0sh512/profhome-vpn-bot}{Marzban Telegram VPN Bot}}{контейнеризированный Python-бот для управления приглашениями и жизненным циклом пользователей через Marzban API.}

\end{document}
